% ==================================================================
% Conscious Point Physics:
% SS-1f — The Physical Realization of the SU(3) Hop
% Strong Sector Series — SU(3) foundations sub-family (1a..1f)
% Companion to: SS-1b (SU(3) algebra, exact), SS-1c (eight gluons),
%               SS-3 (uniqueness), SF-5 (strong-sector flagship)
% Fills: Open Problem op:strong_primitive (registered in SS-1b)
% GitHub: CPP/series_strong/papers/SS-1f_su3_hop_realization.tex
% ------------------------------------------------------------------
% CHANGELOG
%   v0.1  18 Jun 2026 — first assembly (mechanism note). Picture, not
%         derivation. Honesty box front-and-centre. Shared-registry
%         registration DEFERRED to SHIP-time integration (multi-window
%         rule E). Source: su3_mechanism_realization_proposal.md (TODO-019),
%         Session-16x exploratory arc, and the founders_voice mechanism notes.
%   v0.2  18 Jun 2026 — panel cycle (ChatGPT/Grok/Gemini/Copilot). Surgical
%         revisions: (i) Prop 6.1 iff scoped to the vertex-occupancy realisation
%         + Remark on scope; (ii) C^3 has only 2 traceless diagonal d.o.f. (the
%         8-fold torus is for abstract modes only) added to statement+proof;
%         (iii) §1 'physical' -> 'candidate ... pending §7 frame map'; (iv) §7
%         seam sharpened to the explicit 'what does E_ij act on?' sub-problem +
%         labelling-vs-dynamics split (answers Gemini's geometry objection by
%         placing the geometric fixing inside the forcing problem); (v) new
%         Open Problem 9.4 — connection to SS-1d dynamics. Verdict split 2 ship
%         / 2 reconcile-seam; edits address the converged math points; the
%         ship-vs-resolve-seam decision is the remaining fork (TLA to decide).
%   v1.0  18 Jun 2026 — SHIP. Panel: 3/4 ship after v0.2 math edits; remaining
%         fork was the §7 frame seam. TLA clarification RESOLVES it: the hTetra
%         is the baryon BONDING FRAMEWORK, not a per-quark cage; up/down quarks
%         are cageless; strange carries its own hTetra cage. §3/§7 rewritten —
%         colour = baryon-framework vertex occupancy; E_ij acts on baryon-vertex
%         occupancy; the operator-domain question is answered, not registered.
%         §7 flags the tension this surfaces with SS-1b's per-quark-cage EXPOSITION
%         ('each quark sits inside a tetrahedral cage'): SS-1b's ALGEBRA is
%         frame-agnostic and unaffected; its physical-picture wording needs
%         harmonising (registered TODO). §1 + §9 updated (frame no longer 'open';
%         residual = SS-1b harmonisation, a corpus task). Final ChatGPT micro-edits
%         folded (explicit abstract-vs-C^3-diagonal parenthetical; 'candidate
%         physical'). Acknowledgements v0.1->v1.0 leftover fixed (Grok). Shared-
%         registry integration still deferred to a flagged patch.
% ==================================================================

\documentclass[11pt,a4paper]{article}
\usepackage[margin=1in]{geometry}
\usepackage[utf8]{inputenc}
\usepackage[T1]{fontenc}
\usepackage{amsmath,amssymb,amsthm}
\usepackage{graphicx}
\usepackage{caption,float,booktabs,parskip,xcolor}
\usepackage[numbers,sort&compress]{natbib}
\usepackage{hyperref}
\hypersetup{colorlinks=true,linkcolor=blue,urlcolor=cyan,citecolor=red}

\newtheorem{theorem}{Theorem}[section]
\newtheorem{proposition}[theorem]{Proposition}
\newtheorem{corollary}[theorem]{Corollary}
\newtheorem{lemma}[theorem]{Lemma}
\newtheorem{remark}[theorem]{Remark}
\newtheorem{openprob}[theorem]{Open Problem}

\newcommand{\SSV}{\mathrm{SSV}}
\newcommand{\hDP}{\mathrm{hDP}}
\newcommand{\qCP}{\mathrm{qCP}}
\newcommand{\eCP}{\mathrm{eCP}}
\newcommand{\qDP}{\mathrm{qDP}}
\newcommand{\eDP}{\mathrm{eDP}}
\newcommand{\ZBW}{\mathrm{ZBW}}
\newcommand{\phig}{\varphi}
\newcommand{\ket}[1]{|#1\rangle}
\newcommand{\bra}[1]{\langle#1|}

\title{\textbf{Conscious Point Physics:\\
The Physical Realization of the SU(3) Hop\\
{\large How a Charge-Driven ZBW Transition Carries the Colour Generator}}\\[6pt]
{\large Strong Sector Series --- SS-1f --- Version 1.0 (mechanism note)}}

\author{Thomas Lee Abshier, ND \and Claude Opus (Anthropic)\\
Hyperphysics Institute\\
\url{https://hyperphysics.com}}

\date{18 June 2026}

\begin{document}
\maketitle

\begin{abstract}
SS-1b proves that the colour algebra $\mathfrak{su}(3)$ emerges \emph{exactly}
from the three colour vertices of the tetrahedral cage, and SS-1c builds the
eight gluons on those vertices; SS-3 shows the realisation is unique within the
operator representation. All three explicitly leave open
(\emph{op:strong\_primitive}, registered in SS-1b) the prior physical question:
\emph{why} is the strong interaction carried by tetrahedral hopping at all?
This note proposes the physical realisation. Colour is identified with
\emph{vertex occupancy} --- which cage vertex a quark's central $\qCP$ is bonded
to. Each colour generator is realised as a charge-driven
zitterbewegung ($\ZBW$) / Stress-Sea-Vector ($\SSV$) transition in which a quark
leaves the superimposed bond at one vertex and enters it at another, and the
\emph{non-commutativity} of two such moves arises physically because each quark's
$\SSV$ gradient is conditioned by the occupancy of the other vertices.
We isolate the load-bearing constraint as a proposition: the eight
\emph{uncoupled} single-bond oscillation modes generate the \emph{abelian}
algebra $\mathfrak{u}(1)^{\oplus 8}$ (a torus), \emph{not} $\mathfrak{su}(3)$;
the $\mathfrak{su}(3)$ structure constants appear if and only if the generators
are realised as inter-vertex hop operators. This both rules out the naive
``eight modes $=$ eight gluons'' identification and shows precisely what the
mechanism must supply. We are explicit about status: this is a \emph{mechanism
note, not a derivation}. The picture is a \emph{candidate carrier} of the
SS-1b algebra, not a proof that the substrate geometry \emph{forces} it; the
per-quark colour-cage frame of SS-1b and the baryon-$\hDP$-tetrahedron frame
of the founder's mechanism account are not yet reconciled; and the mechanism
is silent inside a stable (energy-eigenstate) baryon, so its distinctive
physics, if any, lives in non-eigenstate processes. The forcing question of
\emph{op:strong\_primitive} therefore remains open.
\end{abstract}

\section{What this note does and does not claim}
\label{sec:honesty}

This section is placed first by design, so that no reader --- including its
authors --- mistakes a picture for a proof.

\medskip
\noindent\textbf{This note \emph{does}:}
\begin{itemize}
\item supply a \emph{candidate physical}, corpus-consistent realisation of the SU(3)
colour hop --- colour as occupancy of the baryon bonding tetrahedron's vertices,
the gluon generator as a charge-driven $\ZBW$/$\SSV$ transition, and the
non-commutativity as a consequence of the coupling between vertices
(Sections~\ref{sec:occupancy}--\ref{sec:coupling}); it remains a \emph{candidate}
because the forcing is open (Section~\ref{sec:status}), not because the frame is
ambiguous (the frame is fixed in Section~\ref{sec:seam});
\item prove a negative structural result --- that uncoupled oscillation modes
form a torus and \emph{cannot} be the generators --- which forces the hop
picture and rules out the naive mode-count identification
(Proposition~\ref{prop:torus}, Section~\ref{sec:torus}).
\end{itemize}

\medskip
\noindent\textbf{This note \emph{does not}:}
\begin{itemize}
\item derive $\mathfrak{su}(3)$. That is SS-1b's theorem~\citep{ss1b}, cited
here and not re-proved. Given colour-as-vertex, the algebra follows from the
three-ness of the vertex set; this note adds the \emph{physical carrier}, not
the algebra;
\item \emph{force} the realisation from substrate geometry. The $\ZBW$/$\SSV$
mechanism is shown to be a candidate \emph{rich enough} to host the generators;
it is not shown that the geometry permits no other carrier. Closing this gap is
closing \emph{op:strong\_primitive};
\item provide a falsifiable in-baryon prediction. A stable baryon is an
energy eigenstate; the hops are real as structure but silent as activity
(Section~\ref{sec:testable}).
\end{itemize}

\medskip
\noindent\textbf{The open items are registered honestly} and collected in
Section~\ref{sec:status}: (i) candidate-carrier vs forcing; (ii) the absence of
an in-baryon test; (iii) consistency with the SS-1d dynamics. The frame question
(per-quark cage vs baryon bonding tetrahedron) is \emph{not} left open: it is
resolved in Section~\ref{sec:seam} in favour of the baryon bonding tetrahedron,
with the harmonisation of SS-1b's expository cage language registered as a corpus
task.

\section{The open problem SS-1b left}
\label{sec:openprob}

SS-1b~\citep{ss1b} establishes, as an exact algebraic identity, that eight
DI-bit hopping operators on the three tetrahedral colour vertices --- six
colour-changing (edge) operators and two diagonal operators --- satisfy
\begin{equation}
T^a = \tfrac{1}{2}\lambda^a,\qquad [T^a,T^b] = i f^{abc} T^c,
\label{eq:su3}
\end{equation}
with the standard Gell-Mann structure constants $f^{abc}$. SS-1c~\citep{ss1c}
builds the eight gluons as the corresponding $\hDP$ configurations (six
colour-changing, two colour-neutral), and SS-3~\citep{ss3} shows the
realisation is unique within the operator representation.

The construction of SS-1b takes as an assumption (its hypothesis~(a)) that the
three colour states \emph{are} the three base vertices, and the hops are moves
among them. It records, as Open Problem~\emph{op:strong\_primitive}, the
physical question this leaves untouched: \emph{why} is the strong force
mediated by tetrahedral hopping among those vertices, rather than by some other
subgraph or some other dynamics entirely? Equation~\eqref{eq:su3} is the
\emph{grammar} of three labelled corners; it does not by itself say what
\emph{physically} performs a move between corners, nor why that performance is
the strong interaction. This note addresses precisely that layer.

\section{Colour as vertex occupancy}
\label{sec:occupancy}

We adopt the identification already implicit in SS-1c's baryon
construction~\citep{ss1c}: a baryon places one quark at each of three vertices
of the baryon bonding tetrahedron, the colour index is the vertex, and the
colour singlet is the totally antisymmetric combination over the three vertices.
The bonding tetrahedron is the framework that binds the quarks (Section~\ref{sec:seam}),
not a per-quark cage --- the up and down quarks are cageless --- so the vertex
in question is a vertex of the shared baryon framework.

\begin{equation}
\text{colour of a quark} \;\equiv\; \text{the baryon-bonding-tetrahedron vertex its central }\qCP
\text{ occupies.}
\end{equation}

Colour is thus not a property painted on the quark in isolation; it is a
\emph{relationship to the frame} --- which seat the quark occupies. The colour
Hilbert space $\mathcal{H}_{\rm colour}\cong\mathbb{C}^3$ is spanned by the
three vertex-occupancy states $\ket{V_1},\ket{V_2},\ket{V_3}$. This is the
fundamental representation, made concrete: the ``3'' of SU(3) is the three-ness
of the base, and the ``8'' is forced from it as $3^2-1$.

\section{The hop as a ZBW / SSV transition}
\label{sec:hop}

Each quark's bond to its vertex is a $\ZBW$ oscillation between two physically
distinct states, following the founder's mechanical account:
\begin{itemize}
\item \textbf{State~1 (away).} The quark sits off the vertex; its $\SSV$ gradient
is dominated by attraction toward the oppositely charged vertex, drawing it in.
\item \textbf{State~2 (superimposed).} The quark sits superimposed on the
vertex; its $\SSV$ gradient is now dominated by the surrounding Dipole Sea,
pushing it back out.
\end{itemize}
The bond is the limit cycle between these: in--out, State~1$\,\leftrightarrow\,$State~2.

A \emph{hop} --- the physical content of a colour generator --- is the event in
which a quark, rebounding from State~2 at vertex $V_j$, re-enters State~2 at a
\emph{different} vertex $V_i$ under the prevailing multi-body $\SSV$ gradient.
Write this transition as the operator $E_{ij}$, ``occupancy moves $V_j\!\to\!
V_i$.'' The $\ZBW$ oscillation is the physical muscle that performs the move; the
hop is the move itself. The in/out asymmetry of the $\ZBW$ bond (the inbound
$\SSV$ profile differs from the outbound one) makes each carried generator
directional, $E_{ij}\neq E_{ji}$, but --- as the next two sections make precise
--- directionality is a property of the \emph{carrier}, not the source of the
algebra.

\section{Non-commutativity from inter-vertex coupling}
\label{sec:coupling}

The defining feature of a non-abelian gauge algebra is that two generators fail
to commute, and that the failure is itself a generator. Physically, this
requires that the action of one move alter the conditions for another. The
mechanism supplies exactly this coupling: \emph{each quark's $\SSV$ gradient is
conditioned by the occupancy of the other vertices.} The gradient that decides
whether and where the quark at $V_1$ may hop depends on which quark sits at
$V_2$ and $V_3$; performing a hop at one vertex therefore changes the conditions
for a hop at another.

Realised as occupancy operators on $\mathcal{H}_{\rm colour}$, the hops obey the
transposition grammar of three labelled states. For the colour-changing
(off-diagonal) hops,
\begin{equation}
[E_{12},\,E_{23}] = E_{13}\quad(\text{up to sign/phase}),
\label{eq:closure}
\end{equation}
and the difference of the two orderings is itself a third hop in the set:
performing $V_2\!\to\!V_1$ then $V_3\!\to\!V_2$ does not equal the reverse, and
the discrepancy is the $V_3\!\to\!V_1$ hop. Forming the Hermitian combinations
$\tfrac{1}{2}(E_{ij}+E_{ji})$ and $\tfrac{1}{2i}(E_{ij}-E_{ji})$ together with
the two independent diagonal (Cartan) differences --- two, not three, because
the vertex occupancies sum to one quark each, which removes the overall
$\mathrm{U}(1)$ --- reproduces the eight generators of
Equation~\eqref{eq:su3} and the SS-1b closure. The inter-vertex coupling is
the \emph{physical origin} of the structure constants.

\section{Why raw oscillation modes are not the generators (the torus result)}
\label{sec:torus}

It is tempting to identify the eight gluons directly with eight $\ZBW$
oscillation modes --- to ``count to eight'' among the bonds and declare the
algebra found. The following proposition shows this identification fails, and in
failing it tells us precisely what the realisation must supply.

\begin{proposition}[Uncoupled modes form a torus; the colour-changing generators must be hops]
\label{prop:torus}
Take the eight putative gluon degrees of freedom to be uncoupled oscillation
amplitudes --- abstract, mutually commuting phases. As such they generate the
abelian algebra $\mathfrak{u}(1)^{\oplus 8}$, a torus, not $\mathfrak{su}(3)$
(this is an abstract uncoupled-mode torus, \emph{not} eight independent diagonal
operators on $\mathbb{C}^3$, which admits at most three).
Moreover, realised \emph{on the colour space} $\mathcal{H}_{\rm
colour}=\mathbb{C}^3$, the diagonal (Cartan) subalgebra is at most
two-dimensional after removing the trace; so at most two of the eight
$\mathfrak{su}(3)$ generators can be diagonal amplitudes, and the remaining six
colour-changing generators \emph{must} act off-diagonally. Consequently, within
this vertex-occupancy realisation, the $\mathfrak{su}(3)$ structure constants of
Equation~\eqref{eq:su3} are reproduced if and only if the six non-Cartan
generators are realised as inter-vertex transposition (hop) operators $E_{ij}$
on $\mathbb{C}^3$.
\end{proposition}

\begin{remark}[Scope of the ``iff'']
\label{rem:iff}
The biconditional is a statement \emph{about the vertex-occupancy realisation},
not a basis-independent uniqueness claim: $\mathfrak{su}(3)$ admits many bases,
and one is free to rotate generators. The content is that in \emph{any}
realisation that identifies colour with occupancy of three vertices, the
colour-changing generators are necessarily off-diagonal, and the eight-fold
abelian (torus) count is available only for abstract modes, never for
operators acting faithfully on the three-dimensional colour space.
\end{remark}

\begin{proof}[Proof sketch]
An oscillation amplitude is diagonal in the occupancy basis; two diagonal
operators commute, so any set of such amplitudes generates an abelian (toral)
algebra, regardless of their number or of any in/out asymmetry of the individual
bonds. On $\mathbb{C}^3$, however, the space of traceless diagonal operators is
only two-dimensional, so the eight-fold abelian count is available only for
abstract modes; realised on the colour space, at most two generators may be
diagonal and the rest are forced off-diagonal. Directionality ($E_{ij}\neq E_{ji}$) is a property of a single operator
relative to its inverse, not a failure of two distinct operators to commute, and
contributes nothing to the bracket. A nonzero bracket $[A,B]$ requires $A$ and
$B$ to act non-diagonally on a shared subspace --- i.e.\ to \emph{permute} basis
states. The traceless Hermitian operators that permute and rephase three basis
states are exactly the $\mathfrak{su}(3)$ generators, and their closure is
Equation~\eqref{eq:su3} (SS-1b). Hence the off-diagonal hop structure is both
necessary and sufficient for the algebra; the diagonal/amplitude content alone
is necessarily abelian.
\end{proof}

\begin{remark}
Proposition~\ref{prop:torus} is the load-bearing result of this note. It
disqualifies any ``$8$ modes $=$ $8$ gluons'' shortcut, including
flavour-indexed or charge-indexed mode counts (which moreover fail
flavour-blindness: e.g.\ an all-same-charge baryon such as $\Delta^{++}$ or
$\Omega^-$ carries the same colour octet, but a count built from
opposite-charge bond pairs does not reproduce it). It also identifies the
mechanism's job exactly: not to produce eight oscillations, but to realise the
six \emph{off-diagonal} inter-vertex hops as physical, charge-driven
transitions --- which Sections~\ref{sec:hop}--\ref{sec:coupling} do, as a
picture.
\end{remark}

\section{The frame: the baryon bonding tetrahedron}
\label{sec:seam}

The physical frame is fixed, and it is \emph{not} a per-quark cage. The
$\hDP$-tetrahedron (hTetra) is the \emph{bonding framework} that binds the quarks
of a baryon: a hybrid tetrahedron built from type-A and type-B $\hDP$s, three of
whose vertices are occupied by quarks and whose fourth is an open $\eCP$ vertex.
The up and down quarks are \emph{cageless} --- each is a $\qCP$ core with a
radial $\ZBW$ $\eCP$ and a polarised cloud, carrying no tetrahedral cage of its
own; it simply bonds to a vertex of the shared baryon framework. (The strange
quark does carry its own hTetra cage, a separate structure relevant to strange
hadrons; it plays no role in the up/down baryon core treated here.) Consequently
the colour of a quark is the baryon-framework vertex its $\qCP$ occupies, and the
hop operator $E_{ij}$ acts on \emph{baryon-framework vertex occupancy}: it moves
a quark between vertices of the single, shared tetrahedron. There is no competing
``internal colour coordinate'' frame for the up/down baryon core, because there
is no per-quark cage. The operator-domain question raised in review --- \emph{what
does $E_{ij}$ act on?} --- is therefore answered: baryon-framework vertex
occupancy.

\paragraph{Relation to SS-1b, and a registered harmonisation.} SS-1b's algebra
is \emph{frame-agnostic}: it requires only three labelled colour vertices and
DI-bit hops among them, and is unaffected by which physical structure carries
those vertices. Its \emph{exposition}, however, describes the three colour
vertices as the base of a per-quark ``$\qCP$ tetrahedral cage'' (apex $V_4$ = that
quark's central $\qCP$), with the phrase ``each quark sits inside a tetrahedral
cage.'' Read literally for an up or down quark, that per-quark-cage language is
in tension with the cageless up/down quarks of the bonding-framework realisation
adopted here. We resolve the frame in favour of the baryon bonding tetrahedron,
and register two harmonisation tasks as corpus items (todolist): (i) reconcile
SS-1b's expository per-quark-cage language with the cageless-quark /
bonding-framework realisation; (ii) work out the detailed correspondence between
SS-1b's (apex-$\qCP$ + three-colour-base) cage description and the baryon hTetra's
(three-quark + open-$\eCP$-vertex) geometry. \textbf{The algebra of SS-1b is not at
stake --- only its physical-picture wording.}

\paragraph{What remains open.} Fixing the frame does not close the forcing
problem. The precise mechanics of a colour-changing hop on a \emph{singly
occupied} baryon tetrahedron --- how the substrate dynamics move a quark's
colour between vertices that are each already occupied --- is part of the forcing
problem (Section~\ref{sec:status}), not of the frame question, which is now
settled.

\section{Where the testable physics lives}
\label{sec:testable}

A stable proton or neutron is an energy eigenstate. Its colour structure is
stationary; no hop is actually firing, and the $\ZBW$ trembling, though real, is
silent as activity. Consequently the realisation \emph{cannot} be tested inside
a stable baryon: there is no live transition to catch. A session-level search
for in-baryon observables (extra short-lived configurations from minor bonding
modes) found them parametrically suppressed.

The mechanism's distinctive physics, if any, therefore lives in
\emph{non-eigenstate} processes --- where the cage is forced out of its ground
configuration: beta decay, meson formation, and high-energy collisions. The
companion beta-decay and chirality mechanism (founder's notes; to be developed)
is where a charge-driven transition becomes an actual event with a direction and
a handedness, and is the natural place to look for a deviation from, or a
substrate-derivation of, the established phenomenology. This note makes no such
prediction; it only locates where one could live.

\section{Status and open problems}
\label{sec:status}

\begin{openprob}[Forcing of the tetrahedral hop --- op:strong\_primitive, mechanism side]
Show that the substrate dynamics \emph{force} the SU(3) colour hop to be the
$\ZBW$/$\SSV$ inter-vertex transition of Sections~\ref{sec:hop}--\ref{sec:coupling},
rather than merely permitting it; equivalently, derive
Equation~\eqref{eq:su3}'s realisation from the substrate rather than adopting
colour-as-vertex as a hypothesis. \textbf{Open.}
\end{openprob}

\begin{openprob}[SS-1b exposition harmonisation (residual of the now-resolved frame)]
The physical frame is resolved (Section~\ref{sec:seam}): the baryon bonding
tetrahedron, with cageless up/down quarks. What remains is a corpus-harmonisation
task, not a physics ambiguity: reconcile SS-1b's expository per-quark
``$\qCP$ tetrahedral cage'' language with the cageless-quark realisation, and
work out the detailed correspondence between SS-1b's (apex-$\qCP$ +
three-colour-base) cage and the baryon hTetra's (three-quark + open-$\eCP$-vertex)
geometry. SS-1b's \emph{algebra} is frame-agnostic and unaffected.
\textbf{Open (corpus harmonisation).}
\end{openprob}

\begin{openprob}[An observable]
Determine whether the realisation predicts any deviation from established
phenomenology in a non-eigenstate process (Section~\ref{sec:testable}), or
prove it cannot. \textbf{Open.}
\end{openprob}

\begin{openprob}[Connection to the strong-sector dynamics]
Connect this kinematic hop realisation to the dynamical strong-sector results
already in the corpus --- the running coupling, confinement, and positive beta
function of SS-1d --- and show that the realisation reproduces, or at least
coexists consistently with, them. This note establishes a kinematic/structural
realisation of the colour algebra only; it does not yet demonstrate that the hop
dynamics yields those dynamical features. \textbf{Open consistency check.}
\end{openprob}

\paragraph{Summary of standing.} The algebra (SS-1b) is proven and is not
touched here. The realisation offered in this note is a \emph{candidate
carrier}, internally consistent and corpus-consistent, supported by the negative
result Proposition~\ref{prop:torus}, which fixes what the carrier must be (an
off-diagonal hop, not an amplitude). It is a \emph{mechanism, not a derivation}.
\emph{op:strong\_primitive} remains open on the forcing question.

\section*{Acknowledgements}
This note records the mechanical account developed by T.L.A.\ during the
Session-16x exploratory arc, with the algebra and uniqueness results due to the
prior SS-1b/SS-1c/SS-3 work. It is issued at v1.0 as a mechanism note;
shared-registry registration is deferred to ship-time integration.

\begin{thebibliography}{9}
\bibitem{ss1b} T.~L.~Abshier, \emph{SU(3)$_c$ Algebra from Tetrahedral Phase
Interference: The Strong Sector Algebra} (SS-1b), Conscious Point Physics,
Hyperphysics Institute (2026).
\bibitem{ss1c} T.~L.~Abshier, \emph{The Eight Gluons as hDP Structures} (SS-1c),
Conscious Point Physics, Hyperphysics Institute (2026).
\bibitem{ss1d} T.~L.~Abshier, \emph{Confinement and the Beta Function} (SS-1d),
Conscious Point Physics, Hyperphysics Institute (2026).
\bibitem{ss3} T.~L.~Abshier, \emph{SU(3) Uniqueness} (SS-3), Conscious Point
Physics, Hyperphysics Institute (2026).
\bibitem{sf5} T.~L.~Abshier and Claude Opus, \emph{Strong-Sector Unification
from 600-Cell Geometry} (SF-5), Conscious Point Physics Flagship Series,
Hyperphysics Institute (2026).
\end{thebibliography}

\end{document}
